\appendixsection{线性表示上的运算}


设 \(G\) 是一个群，\(k\) 是交换域，\(E_1, E_2, \ldots, E_n\) 是 \(k\) 上的向量空间，\(\pi_i\) 是 \(G\) 在 \(E_i\) 上的线性表示（\(1 \leqslant i \leqslant n\)）。映射

\[
g \mapsto \pi_ {1} (g) \otimes \dots \otimes \pi_ {n} (g)
\]

是从 G 到向量空间 \(E_{1} \otimes \cdots \otimes E_{n}\) 的线性映射，称为 \(\pi_{1}, \ldots, \pi_{n}\) 的张量积，记作 \(\pi_{1} \otimes \cdots \otimes \pi_{n}\)。

设 E 是 k 上的向量空间，\(\pi\) 是 G 在 E 上的线性表示。对所有 \(g \in G\)，令 \(\tau(g)\)（分别地，\(\sigma(g)\)、\(\varepsilon(g)\)）为代数 \(\mathsf{T}(\mathsf{E})\)（分别地，\(\mathsf{S}(\mathsf{E})\)、\(\bigwedge(\mathsf{E})\)）中延拓 \(\pi(g)\) 的唯一自同构（《代数》第三章第 5 节第 2 号、第 6 节第 2 号和第 7 节第 2 号）。于是 \(\tau\)（分别地，\(\sigma\)、\(\varepsilon\)）是 G 在 T(E)（分别地，\(\mathsf{S}(\mathsf{E})\)、\(\bigwedge(\mathsf{E})\)）上的线性表示，记作 \(\mathsf{T}(\pi)\)（分别地，\(\mathsf{S}(\pi)\)、\(\bigwedge(\pi)\)）。由 T(π)（分别地，\(\mathsf{S}(\pi)\)、\(\bigwedge(\pi)\)）定义于 \(\mathsf{T}^{\mathsf{n}}(\mathsf{E})\)（分别地，\(\mathsf{S}^{\mathsf{n}}(\mathsf{E})\)、\(\bigwedge^{\mathsf{n}}(\mathsf{E})\)）的子表示，称为 \(\pi\) 的 n 次张量（分别地，对称、外）幂，记作 \(\mathsf{T}^{\mathsf{n}}(\pi)\)（分别地，\(\mathsf{S}^{\mathsf{n}}(\pi)\)、\(\bigwedge^{\mathsf{n}}(\pi)\)）。于是

\[
T ^ {n} (\pi) = \pi \otimes \pi \otimes \dots \otimes \pi
\]

（n 个因子）。表示 \(\mathsf{S}(\pi), \bigwedge (\pi)\) 是 \(\mathsf{T}(\pi)\) 的商表示，因此 \(\mathsf{S}^n (\pi), \bigwedge^n (\pi)\) 是 \(\mathsf{T}^n (\pi)\) 的商表示。

设 g 是 k 上的李代数。有限个 g 的表示的张量积已在第一章第 3 节第 2 号定义；记作 \(\pi_{1}\otimes\cdots\otimes\pi_{n}\)。设 E 是 k 上的向量空间，\(\pi\) 是 g 在 E 上的表示。对所有 \(x\in g\)，令 \(\tau'(x)\)（分别地，\(\sigma'(x)\)、\(\varepsilon'(x)\)）为代数 \(\mathrm{T(E)}\)（分别地，\(\mathrm{S(E)}\)、\(\bigwedge(\mathrm{E})\)）中延拓 \(\pi(x)\) 的唯一导子（《代数》第三章第 10 节第 9 号例 1）。于是由《代数》同上，公式 (35)，\(\tau'\)（分别地，\(\sigma'\)、\(\varepsilon'\)）是 g 在 T(E)（分别地，S(E)、\(\bigwedge(\mathrm{E})\)）上的线性表示，

记作 \(\mathsf{T}(\pi)\)（分别地，\(\mathsf{S}(\pi),\bigwedge (\pi)\)）。由 \(\mathsf{T}(\pi)\)（分别地，\(\mathsf{S}(\pi),\bigwedge (\pi)\)）定义的、由 \(\mathsf{T}^n (\mathsf{E})\)（分别地，\(\mathsf{S}^n (\mathsf{E}),\bigwedge^n (\mathsf{E}))\)）决定的子表示，记作 \(\mathsf{T}^n (\pi)\)（分别地，\(\mathsf{S}^n (\pi),\bigwedge^n (\pi))\)）。表示 \(\mathsf{T}^n (\pi)\) 是 \(n\) 个与 \(\pi\) 相同的表示的张量积。表示 \(\mathsf{S}(\pi),\bigwedge (\pi)\) 是 \(\mathsf{T}(\pi)\) 的商表示，因此 \(\mathsf{S}^n (\pi),\bigwedge^n (\pi)\) 是 \(\mathsf{T}^n (\pi)\) 的商表示。

